\documentclass[11pt]{article}
\usepackage[margin=1in]{geometry}
\usepackage{setspace}
\usepackage{enumitem}
\usepackage{amsmath}
\usepackage{amssymb}
\usepackage{fancyhdr}
\IfFileExists{lastpage.sty}{\usepackage{lastpage}}{}
\usepackage{hyperref}

\setlength{\parindent}{0pt}
\setlength{\parskip}{0.6em}

\pagestyle{fancy}
\fancyhf{}
\IfFileExists{lastpage.sty}{
  \fancyfoot[C]{Financial Data Science II\ \ \ Updated March 10, 2026\ \ \ Page \thepage\ of \pageref{LastPage}}
}{
  \fancyfoot[C]{Financial Data Science II\ \ \ Updated March 10, 2026\ \ \ Page \thepage}
}

\begin{document}

\begin{center}
{\Large \textbf{Part 22: Feature Engineering}}
\end{center}

\textbf{Part 22: Feature Engineering}

Feature engineering is the process of transforming raw data into signals
that are informative for a prediction or decision task. In financial ML,
this often makes the difference between a backtest that works and a model
that survives in production.

The focus in this note is on feature construction, validation, and
organization of a feature library for predictive modeling.

\section*{Core Principles}

\begin{itemize}[leftmargin=*]
  \item \textbf{Causality and timing:} Features must be computable at the
  prediction time. Any accidental use of future information creates
  lookahead bias.
  \item \textbf{Stability:} Favor features whose relationship to the target
  is stable across time and regimes.
  \item \textbf{Interpretability:} When possible, maintain financial
  intuition for each feature family.
  \item \textbf{Cost awareness:} Favor features that can be computed cheaply
  and consistently at scale.
\end{itemize}

\section*{Feature Families}

\textbf{1. Returns and Momentum}

\begin{itemize}[leftmargin=*]
  \item Simple and log returns at multiple horizons
  \item Momentum and risk-adjusted momentum
  \item Acceleration of returns or momentum
\end{itemize}

\textbf{2. Volatility and Risk}

\begin{itemize}[leftmargin=*]
  \item Rolling volatility (standard and exponential)
  \item Volatility of volatility
  \item Drawdown and time since peak
  \item ATR and normalized range
\end{itemize}

\textbf{3. Trend and Technical Indicators}

\begin{itemize}[leftmargin=*]
  \item Moving averages and ratios
  \item RSI, MACD, Bollinger features
  \item Trend strength and regime proxies
\end{itemize}

\textbf{4. Liquidity and Volume}

\begin{itemize}[leftmargin=*]
  \item Volume z-scores and volume ratios
  \item Dollar volume and liquidity proxies
  \item OBV / VPT style cumulative volume signals
\end{itemize}

\textbf{5. Cross-Asset Features}

\begin{itemize}[leftmargin=*]
  \item Rolling correlation with benchmark
  \item Beta and alpha versus benchmark
  \item Relative momentum vs. index
\end{itemize}

\section*{Predictive Feature Workflow}

\begin{enumerate}[leftmargin=*]
  \item \textbf{Define target:} Choose a forward return horizon
  (e.g., 5d, 20d). Decide regression or directional classification.
  \item \textbf{Construct features:} Build feature groups above.
  \item \textbf{Temporal validation:} Walk-forward evaluation and
  leakage checks. Avoid sample re-use.
  \item \textbf{Feature screening:} Rank features using linear and nonlinear
  tests (e.g., F-test and mutual information).
  \item \textbf{Ablation:} Add feature families one at a time to show
  incremental out-of-sample value.
\end{enumerate}

\section*{Exercises}

\textbf{Exercise 1.} Build a small feature set with returns, volatility,
RSI, and drawdown. Use a 5-day forward return target and compare
walk-forward performance for each feature family.

\textbf{Exercise 2.} Add a cross-asset benchmark feature (beta vs. SPY)
for the same target. Measure whether it improves out-of-sample MSE or
hit rate.

\section*{Notes and Code References}

A working feature library and example scripts are in:

\begin{itemize}[leftmargin=*]
  \item \texttt{/Users/harold/4. RA work/yfinance\_feature\_engineering}
  \item \texttt{/Users/harold/4. RA work/FACTOR\_TRAINING\_CHI}
\end{itemize}

These repositories include implementations of the feature families
listed above, along with walk-forward validation and feature ranking.

\end{document}
